\section{Etwas Mathematik}
\label{sec:etwas-mathematik}

%%%%%%%%%%%%%%%%%%%%%%%%%%%%%%%%%%%%%%%%%%%%%%%%%%%%%%%%%%%%%%%%%%%%%%%%%%%%%%%
%%%%%%%%%%%%%%%%%%%%%%%%%%%%%%%%%%%%%%%%%%%%%%%%%%%%%%%%%%%%%%%%%%%%%%%%%%%%%%%
%%%%%%%%%%%%%%%%%%%%%%%%%%%%%%%%%%%%%%%%%%%%%%%%%%%%%%%%%%%%%%%%%%%%%%%%%%%%%%%
\subsection{Das Binärsystem}
\label{sec:etwas-mathematik:binärsystem}

Computer rechnen im sogenannten Binärsystem, sie kennen nur 0 und 1. Um größere Zahlen darstellen zu können, werden mehrere dieser Nullen und Einsen, genannt ``\textbf{Bit}'' aneinandergereit. Am gebräuchlichsten dabei ist das ``\textbf{Byte}'' (siehe Abb.~\ref{fig:byte_bits}), das aus 8 Bit besteht.

\begin{figure}[h]
  \centering
  \begin{tikzpicture}[
      cell/.style={
          draw,
          rectangle,
          minimum width=0.8cm,
          minimum height=0.8cm,
          font=\ttfamily\large,
          outer sep=0pt,          % Kein Abstand zwischen Zellen
          anchor=base,
      }
  ]

  % Bit-Positionen (0 bis 7) nebeneinander
  \foreach \i in {0,1,...,7} {
      \node[cell] at (7-\i,0) (bit\i) {\i};
  }

  % Border around entire byte
  \node[draw=black, thick, inner sep=4pt, fit=(bit0)(bit7)] (bytebox) {};

  % Etikett "Bit" darüber
  \node[font=\sffamily, left=0.1cm of bytebox] {\textbf{Bit}};

  \end{tikzpicture}
  \caption{Ein Byte mit den Bit-Positionen 0 bis 7.}
  \label{fig:byte_bits}
\end{figure}

\noindent
Jedem Bit wird nun ein Wert zugeordnet. Das erfolgt mathematisch gemäß:
\begin{equation*}
  2^b
\end{equation*}
wobei $b$ die Nummer des Bits ist. Das ist eine sogenannte ``Potenz'', eine mathematische Operation die sagt, wie oft eine Zahl mit sich selbst multipliziert wird. In diesem Fall wird die ``2'' $b$-mal mit sich selbst multipliziert, zum Beispiel:
\begin{align*}
  2^1 &= 2              &= 2\\
  2^2 &= 2\cdot2        &= 4\\
  2^3 &= 2\cdot2\cdot2  &= 8
\end{align*}
Interessant ist nun die Frage, was $2^0$ ist. Wenn ihr genau hinschaut, dann merkt ihr, dass, um in die nächste Zeile zu kommen, ihr das Ergebnis mit $2$ multipliziern müsst. Umgekehrt genauso: um in die vorhergehende Zeile zu kommen, muss man durch $2$ teilen. Wenn also $2^1$ zwei ist und man, um zu $2^0$ zu kommen durch zwei Teilen muss, dann ergibt sich
\begin{equation*}
  2^0 = 1\,.
\end{equation*}
Übrigens, jede Zahl hoch $0$ hat als Ergebnis die $1$ (wir kommen später darauf zurück).

Um auf das \textbf{Byte} zurückzukommen, heißt dass, dass den \textbf{Bits} jeweils die Werte 1 bis 128 zugeordnet sind (siehe Abb.~\ref{fig:byte_values}).

\begin{figure}[h]
  \centering
  \begin{tikzpicture}[
      cell/.style={
          draw,
          rectangle,
          minimum width=0.8cm,
          minimum height=0.8cm,
          font=\ttfamily\large,
          outer sep=0pt,          % Kein Abstand zwischen Zellen
          anchor=base,
      }
  ]

  % Zeile 1: Bit-Positionen (0 bis 7)
  \foreach \i in {0,1,...,7} {
      \node[cell] at (7-\i,0) (bit\i) {\i};
  }

  % Zeile 2: Werte (1, 2, 4, 8, 16, 32, 64, 128)
  \foreach \i in {0,1,...,7} {
      \node[cell] at (7-\i,-1.2) (val\i) {\pgfmathparse{int(2^\i)}\pgfmathresult};
  }

  % Border around
  \node[draw=black, thick, inner sep=4pt, fit=(bit0)(bit7)] (bytebox) {};
  \node[draw=black, thick, inner sep=4pt, fit=(val0)(val7)] (valbox) {};


  % Etiketten "Bit" und "Wert" darüber
  \node[font=\sffamily, left=0.1cm of bytebox] {\textbf{Bit}};
  \node[font=\sffamily, left=0.1cm of valbox] {\textbf{Wert}};

  \end{tikzpicture}
  \caption{Ein Byte mit Bit-Positionen und den zugehörigen Werten ($2^b$).}
  \label{fig:byte_values}
\end{figure}

\noindent
Damit können wir nun die Zahlen 0 bis 255 darstellen, z.B.:
\begin{align*}
  00000000 &\rightarrow 0\,, \\
  00000011 &\rightarrow 3\,, \\
  01011000 &\rightarrow 88\,, \\
  11111111 &\rightarrow 255\,.
\end{align*}
Das Addieren funktioniert nun genauso, wie ihr es aus der Schule, in der natürlich das Zehnersystem verwendet wird, gelernt habt. Beispiel:
\begin{lstlisting}
      00100111       39
    + 00000110     +  6
          11         1     (Übertrag)
     ---------      ---
      00101101       45
\end{lstlisting}
Entsprechend lässt sich auch die Multiplikation durchführen:
\begin{lstlisting}
      00001010·00101010 =     10·42
               ..|.|.|.
        00000000.|.|.|.
         00000000|.|.|.
          00001010.|.|.
           00000000|.|.
            00001010.|.
             00000000|.
              00001010.
               00000000
               1 1             (Übertrag)
        ---------------
              110100100        420
\end{lstlisting}
Wenn ihr genau hinseht, dann bemerkt ihr, dass wir zwei Bytes multipliziert haben, aber das Ergebnis benötigt ein Bit mehr. Das führt zu dem Problem, dass, wenn ihr auf einem Computer diese Multiplikation durchführt, das höchste Bit (die letzte ``1'' links) abgeschnitten wird und das Ergebnis, dass der Computer euch liefert, $164$ ist. Das gleiche passiert bei einem Byte, wenn ihr zu $255$ die $1$ dazuzählt, das Ergebnis ist $0$ (in der Fachsprache als \textbf{Überlauf} bezeichnet).

Computer können auch negative Zahlen darstellen. Dazu nützt man aus, dass, wenn $255 + 1$ zum Ergebnis $0$ hat, dann hat $0-1$ als Ergebnis $-1$, binär also 11111111. Dies nennt man das ``Zweierkomplement'' (den Begriff müsst ihr euch natürlich nicht merken). Wichtig ist: dass höchste Bit zeigt, dass es sich um eine negative Zahl handelt. Damit können dann mit einem Byte die Zahlen -128 bis 127 dargestellt werden:
\begin{lstlisting}
    10000000        -128
    10000001        -127
    ...
    11111111        -  1
    00000000           0
    00000001           1
    ...
    01111111         127
\end{lstlisting}
Wie wird nun aus einer positiven Zahl eine Negative? Das passiert, indem man alle Bits invertiert (das heißt, aus 0 wird 1 und aus 1 eine 0) und 1 addiert:
\begin{lstlisting}
(-127)  10000001  ->  01111110
                     +00000001
                      --------
                      01111111     (127)
\end{lstlisting}
Das funktioniert auch umgekehrt, also aus einer negativen Zahl eine positive zu machen:
\begin{lstlisting}
(127)   01111111  ->  10000000
                     +00000001
                      --------
                      10000001     (-127)
\end{lstlisting}

%%%%%%%%%%%%%%%%%%%%%%%%%%%%%%%%%%%%%%%%%%%%%%%%%%%%%%%%%%%%%%%%%%%%%%%%%%%%%%%
%%%%%%%%%%%%%%%%%%%%%%%%%%%%%%%%%%%%%%%%%%%%%%%%%%%%%%%%%%%%%%%%%%%%%%%%%%%%%%%
%%%%%%%%%%%%%%%%%%%%%%%%%%%%%%%%%%%%%%%%%%%%%%%%%%%%%%%%%%%%%%%%%%%%%%%%%%%%%%%
\subsection{Binäre Operationen — wie Computer rechnen}
\label{sec:etwas-mathematik:binäre-operationen}

Bevor wir in das Thema einsteigen, wie Computer rechnen, sollten wir kurz den 
Begriff ``\textbf{Operation}'' klären.

Eine \textbf{Operation} ist eine \emph{Ausführungshandlung}, also ein Vorgang, 
bei dem aus einem oder mehreren Werten ein neuer Wert berechnet wird. Die 
\textbf{Zeichen}, die diese Operationen beschreiben (wie \texttt{+}, \texttt{-}, 
\texttt{\&}, \texttt{|}), nennt man \textbf{Operatoren}. Die Werte, mit denen 
eine Operation arbeitet, heißen \textbf{Operanden}. Zum Beispiel sind die 
Grundrechenarten, die ihr in der Schule gelernt habt, also Addition, Subtraktion, 
Multiplikation und Division (Plus, Minus, Mal, Geteilt), sogenannte 
\textbf{arithmetische Operationen}. 

Obwohl Computer auch diese arithmetische Operationen durchführen können, arbeiten 
die Schaltungen im Prozessor, dem Herzstück eines Computers, mit sogenannten 
binären Operationen. Auch schon die ersten Computer, die in den 1940er Jahren 
gebaut wurden, arbeiteten mit diesen \emph{binären Operationen}\footnote{%
  Der Autor möchte an dieser Stelle darauf hinweisen, dass diese auch als ``logische
  Operationen'' bezeichnet werden, die Begriffe ``binäre Operation'' und 
  ``logische Operation'' jedoch nicht immer gleich verwendet  werden. Für die Programmiersprache 
  C und C++ gibt es aber eine klare Unterscheidung, auf die wir erst später eingehen werden.} 
(siehe z.B.\ \cite{wikipedia:computer}).

\vskip 1em
\noindent
Die grundlegenden binären Operationen sind:
\begin{description}
  \item[\texttt{AND}]\hfill\\
    Die \texttt{AND}-Operation (binäres UND) verknüpft zwei binäre Zahlen bitweise. Das Ergebnis ist 1, wenn beide Bits 1 sind, sonst ist das Ergebnis 0. Zum Beispiel:
    \begin{lstlisting}
      1101
    & 1011
    --------
      1001
    \end{lstlisting}
  \item[\texttt{OR}]\hfill\\
    Die \texttt{OR}-Operation (binäres ODER) verknüpft zwei binäre Zahlen bitweise. Das Ergebnis ist 1, wenn mindestens eines der Bits 1 ist, sonst ist das Ergebnis 0. Zum Beispiel:
    \begin{lstlisting}
       1101
     | 1011
    --------
       1111
    \end{lstlisting}
  \item[\texttt{XOR}]\hfill\\
    Die \texttt{XOR}-Operation (exklusives ODER) verknüpft zwei binäre Zahlen bitweise. Das Ergebnis ist 1, wenn genau eines der Bits 1 ist, sonst ist das Ergebnis 0. Zum Beispiel:
    \begin{lstlisting}
       1101
     ^ 1011
    --------
       0110
    \end{lstlisting}
  \item[\texttt{NOT}]\hfill\\
    Die \texttt{NOT}-Operation (binäres NICHT) invertiert die Bits einer binären Zahl. Das Ergebnis ist 1, wenn das Bit 0 ist, und 0, wenn das Bit 1 ist. Zum Beispiel:
    \begin{lstlisting}
     ~ 1101
    --------
       0010
    \end{lstlisting}
\end{description}

